
\section{OA-07: Cavity squeezing from a Langevin equation}
\source{2025 PHYS4055 Exam, Question 4(a)--(d); Week 5 open-systems lectures.}
\subsection*{Step 1: obtain the parametric interaction term}
The interaction Hamiltonian and bare cavity equation are
\[
H_I=\frac{i\hbar\chi}{2}(a^2-a^{\dagger2}),\quad
\dot a\big|_{\chi=0}=-\tfrac\gamma2 a+\sqrt\gamma A_{\rm in}.
\]
Using \([a^{\dagger2},a]=-2a^\dagger\),
\[
\frac{i}{\hbar}[H_I,a]
=\frac{i}{\hbar}\left(\frac{i\hbar\chi}{2}
[\,a^2-a^{\dagger2},a]\right)
=-\chi a^\dagger.
\]
Therefore
\[
\boxed{\dot a=-\chi a^\dagger-\tfrac\gamma2a
+\sqrt\gamma A_{\rm in}}.
\]

\subsection*{Step 2: transform the measured quadrature}
Define \(X_+=a+a^\dagger\), and the corresponding
input \(X_{+,\rm in}=A_{\rm in}+A^\dagger_{\rm in}\).
Adding the equation to its adjoint gives
\[
\boxed{\dot X_+=-(\gamma/2+\chi)X_+
+\sqrt\gamma X_{+,\rm in}}.
\]
Take the Fourier convention
\(\tilde X(\omega)=\int_{-\infty}^\infty
e^{-i\omega t}X(t)\,dt\), so
\(\widetilde{\dot X}=i\omega\tilde X\).
Then
\[
\boxed{\tilde X_+(\omega)=
\frac{\sqrt\gamma}{\gamma/2+\chi+i\omega}
\tilde X_{+,\rm in}(\omega)}.
\]

\subsection*{Step 3: apply input--output theory}
With the convention in the exam,
\(A_{\rm out}=\sqrt\gamma\,a-A_{\rm in}\),
so \(X_{+,\rm out}=\sqrt\gamma X_+-X_{+,\rm in}\).
Consequently
\[
\boxed{\tilde X_{+,\rm out}=
\frac{\gamma/2-\chi-i\omega}
{\gamma/2+\chi+i\omega}\,
\tilde X_{+,\rm in}.}
\]
For a vacuum input with unit noise spectral density
and stable dynamics, the output spectral noise is
\[
\boxed{V_+(\omega)=
\frac{(\gamma/2-\chi)^2+\omega^2}
{(\gamma/2+\chi)^2+\omega^2}}.
\]
For \(0<\chi<\gamma/2\), \(V_+<1\) at any
finite frequency, approaching 1 at high frequency.
At \(\omega=0\),
\(V_+(0)=[(\gamma/2-\chi)/(\gamma/2+\chi)]^2\).
The conjugate quadrature is anti-squeezed.
\checkmethod{State the Fourier convention and output-field
sign before doing algebra; otherwise phase factors can appear
inconsistent. A \emph{spectral variance} is based on fluctuation
correlations, not the squared mean field.}
